\scanpage{62}
% archive foliation: 61

\margin{$(\unsure{g_2|g'_2})$}
Cela marche aussi si \deleted{$\unsure{\mathcal E|X}$} $Y$ \inserted{cpt} admet des $\mathrm{rec.}^{\mathrm{ts}}$ arbitrairement fins \inserted{finis} par des $Y_i$ \inserted{fermés ou ouverts} tels que les intersections
$Y_{i_1}\cap Y_{i_2}\cap\cdots\cap Y_{i_\nu}$ $(\nu\ge1)$ de tels $Y_i$ soient $\in\mathcal{E}_{\chi!}$ ....

Dans le cas de $\chi$, on utilise la suite spectrale d'\uline{homologie}, on est ramené à :

\uline{Lemme} Soit $\mathcal R$ syst. local de $k$-vectoriels sur $Y\in\mathcal E_\chi$, fibres de rang \inserted{fini} $b$. Alors
% Two short vertical grouping strokes stand at the left of this lemma.
\figure{\begin{tikzpicture}\draw (0,0)--(0,.65);\draw (0,-.45)--(0,-.2);\end{tikzpicture}}
\[
\chi(Y,\mathcal R)\overset{\mathrm{déf}}=\sum_p(-1)^p
\underbrace{H_{\inserted{p}}^{\deleted{\unsure{p}}}(Y,\mathcal R)}_{\text{le dual est }H^p(Y,\check{\mathcal R})}
\]
existe et :
\[
\chi(Y,\mathcal R)=b\chi(Y).
\]
Cela marche pourvu qu'on fasse l'hyp. \deleted{\unsure{que}} \inserted{\unsure{que}} $Y$ est $\sim Y'$ satisfaisant la condition $(\unsure{g_2|g'_2})$.

\figure{\begin{tikzpicture}\draw (0,0)--(.65,0);\draw (.15,.08)--(.52,.08);\end{tikzpicture}}

\deleted{\unsure{On va se borner}}

\textcircled{h} Définissons $\mathcal C_{-1}=$ \unsure{Ens.|cls.} réduit \unsure{:|à} l'espace vide
et si $n\ge0$ $\mathcal C_n^{\deleted{\unsure{\mathrm{loc}|\mathrm{cpt}}}}=$ \deleted{catégorie} \inserted{classe} des espaces cpts \inserted{$X$} tels que tt pt \inserted{$x\in X$} ait un voisinage $C_x$ tel qu'il existe un espace compact $S\in\unsure{\mathcal E|\mathcal C}_{n-1}$ et un homéom. $\operatorname{Cône}(S)\simeq C_x$ transformant le sommet du cône en $x$.

\[
\mathcal C_\infty=\bigcup_n\mathcal C_n\qquad\text{(espaces loc. côniques)}
\]
\[
\mathcal E_\infty=\deleted{\text{esp}}\text{ des espaces loc. cpts }\inserted{\text{séparés}}X\text{ tels que }\widehat X
\]
soit loc. cônique i.e. $\in C_\infty$.

\margin{\uline{NB}
$\mathcal E_0$ espaces discrets

$\mathcal E_1$ 1-complexes topologiques

$\mathcal E_2$ espaces loc. triangulables de dim $\le2$ [ils sont triangulables dans le cas cpt]

$\mathcal E_3$ espaces loc. triangulables de dim $\le3$ \deleted{[\unsure{ils sont triang. dans} le cas cpt]}

$\mathcal E_4$ esp. loc.\,$\cdots$\,dim $\le4$ (Moïse)

S'ils sont cpts, il n'est pas clair qu'ils sont triangulables ....}
% Right grouping brace beside the first marginal list; angular bracket beside its continuation.
\figure{\begin{tikzpicture}\draw (0,1)..controls(.15,.9)and(.02,.75)..(.18,.6)..controls(.02,.45)and(.15,.2)..(0,0);\draw (.25,0)--(.42,0)--(.42,-.7)--(.25,-.7);\end{tikzpicture}}

\uline{Prop} $\unsure{\mathcal C|\mathcal E}_\infty\subset\mathcal E_\chi\cap\mathcal E_{\chi!}$

Soit $X\in\mathcal E_\infty$. Si $X$ est cpt, \inserted{donc $X\in C_\infty$} comme $X$ est loc. contractile,
on sait que \deleted{\unsure{$X\in$}}$\mathcal E_\chi$, donc $X\in\mathcal E_{\chi!}$. \unsure{Cas|Car} \deleted{\unsure{$X$ général}} \inserted{\unsure{général}} : $X$ loc.
cpt $=\widehat X-(\infty)$, avec $\widehat X\in\mathcal E_\infty$, on a $\widehat X\in\mathcal E_{\chi!}$, donc
$\widehat X-\infty=X\in\mathcal E_{\chi!}$ par (b). Pour prouver $X\in\mathcal E_\chi$, prenons un
$C_x\simeq C(S)$, prenons un homothétique $\frac12 C_x=T$, et

